\documentclass[10pt,a4paper]{amsart}
\usepackage[margin=19mm]{geometry}
\usepackage{amsmath,amssymb,mathtools}
\usepackage[scr=rsfs,cal=euler]{mathalfa}
\usepackage{xfrac}
\usepackage[unicode,hidelinks,bookmarks=false]{hyperref}
\newcommand{\ie}{i.e.\ }
\newcommand{\cf}{cf.\ }
\newcommand{\resp}{respectively\ }
\newcommand{\Subsection}{\subsection}
\providecommand{\nobreakdash}{\nobreak\hbox{-}}
\setlength{\emergencystretch}{2em}
\multlinegap0pt
\usepackage{amssymb}
\usepackage{tikz}
\usepackage{enumerate}
\newtheorem{theorem}{Theorem}
\newtheorem{proposition}{Proposition}
\newtheorem{lemma}{Lemma}
\title{$L$-functions of elliptic curves in ring class extensions of real quadratic fields via regularized theta liftings}
\author{Jeanine Van Order}
\date{Source: arXiv:2510.10277v2 (CC BY 4.0)}
\begin{document}
\maketitle

\noindent\emph{Source and licence.} $L$-functions of elliptic curves in ring class extensions of real quadratic fields via regularized theta liftings. Version arXiv:2510.10277v2 (CC BY 4.0), distributed by arXiv under the Creative Commons Attribution 4.0 International licence, http://creativecommons.org/licenses/by/4.0/, as recorded in the arXiv licence metadata for this exact version and shown on the abstract page https://arxiv.org/abs/2510.10277v2. Full licence notice: \texttt{licenses/arxiv-2510.10277-CC-BY-4.0.txt}.\par\smallskip

\noindent\textit{Editorial scope.} Counted unit: the article's proposition vanish (source0.tex lines 1440--1449). The article's complete original proof of it is delivered with it (source0.tex lines 1451--1479, one \texttt{proof} environment of 29 lines). No mathematics is written, repaired or re-derived: the statement and the proof are the article's own lines, in the article's own notation, and no retained line is substituted or re-flowed.  Retained uncounted as the source-local context the proof needs: lines 1289--1294; lines 1341--1345; lines 1396--1397; lines 1400--1404; lines 1428--1429; lines 1434--1435.  Nothing was omitted from any retained span, so the package carries no omission record.  No retained line contains a citation, so the wrapper carries no bibliography and no bibliography entry.  No retained line contains a figure environment, an image-inclusion command or drawing code, so no image is delivered and the package declares no asset dependency.  Editorial formatting only, in addition: an AMS \texttt{amsart} preamble that restates the article's own declarations and macros used by the retained lines, namely the environments \texttt{theorem}, \texttt{proposition}, \texttt{lemma} on separate counters, together with the standard packages those lines need.  The retained environments print in this document as Theorem 1, Proposition 1, Lemma 1 in order of appearance, whereas the article prints the same items with the numbering of its own full document.  The complete native source of the article as distributed in the arXiv e-print for this exact version is bundled unchanged as \texttt{sources/186873/source0.tex}.
\begin{theorem}[Siegel-Weil for $\mathcal{S}_{L_j}$-valued forms]\label{SW-vector} 

We have the identification of functions of $\tau \in \mathfrak{H}$:
\begin{align*} \frac{\kappa}{2} \cdot \int_{ \operatorname{SO}(V_j)({\bf{Q}}) \backslash \operatorname{SO}(V_j)({\bf{A}})}
\theta_{L_j}(\tau, z, h_f) dh &= E_{L_j}(\tau, s_0, k) = E_{L_j}(\tau, s_0(V_j); k(V_j)). \end{align*}
Here again, $s_0 = s_0(V_j) := \dim(V_j)/2 -1$, and $k = k(V_j) := (p(V_j) - q(V_j))/2 = 0$. \end{theorem}

\begin{proposition}\label{MSEis} 

The Eisenstein series $E_{L_{A,2}}^{\star}(\tau, s)$ has a meromorphic 
continuation to all $s \in {\bf{C}}$, and satisfies the symmetric functional equation 
$E^{\star}_{L_{A,2}}(\tau, s) = E_{L_{A,2}}^{\star}(\tau, -s).$ \end{proposition}

\begin{equation}\begin{aligned}\label{2.20} 
L_2 E_{L_2}(\tau, s; 2) &= \frac{1}{2} \cdot (s-1) \cdot E_{L_2}(\tau, s; 0). \end{aligned}\end{equation}

\begin{proposition}\label{incoherent} 

The nonholomorphic weight-two Eisenstein series $E_{L_2}(\tau, s; 2)$ has a meromorphic continuation 
$E^{\star}_{L_2}(\tau, s; 2) := \Lambda(s+1, \eta) E_{L_2}(\tau, s; 2)$ to all $s \in {\bf{C}}$,
and satisfies the odd functional equation \begin{align*} E_{L_2}^{\star}(\tau, s; 2) = - E_{L_2}^{\star}(\tau, -s; 2).\end{align*} \end{proposition}

\begin{equation}\begin{aligned}\label{2.21}  2 L_2 E_{L_2}'(\tau, 0; 2) &=
E_{L_2}(\tau, 0; 0) - E_{L_2}'(\tau, 0; 0). \end{aligned}\end{equation}

\begin{lemma}\label{diff} The weight-lowering operator $L_l$ can be described in terms of differential forms as 
\begin{align*} \overline{\partial}(f d \tau) &= -v^{2-l} \overline{\xi_l(f)} d \mu(\tau) = - L_l f d \mu (\tau). \end{align*} \end{lemma}

\begin{proposition}\label{vanish}

We have that $E^{\star \prime}_{L_2}(\tau, 0;0) =0$, and hence via $(\ref{2.21})$ that 
$- 2 L_2 E_{L_2}'(\tau, 0; 2) = - E_{L_2}(\tau, 0; 0)$.
Expressed equivalently in terms of differential forms via Lemma \ref{diff}, we obtain the relation
\begin{align*} - 2 L_2 E'_{L_2}(\tau,0; 2) d \mu(\tau) &= 2 \overline{\partial} \left( E'_{L_2}(\tau, 0; 2) d \tau \right)
= - E_{L_2}(\tau, 0; 0) d\mu(\tau), \end{align*} equivalently 
\begin{equation}\begin{aligned}\label{2.23} E_{L_2}(\tau, 0; 0) d \mu(\tau) 
&= - 2 \overline{\partial} \left( E_{L_2}'(\tau, 0; 2) d \tau \right) . \end{aligned} \end{equation}
\end{proposition}

\begin{proof}

We know by the Siegel-Weil formula (Theorem \ref{SW-vector}) 
that the Eisenstein series $E_{L_2}(\tau, s; 0)$ is analytic at $s=0$. Hence, $E_{L_2}(\tau, s; 0)$ 
and its derivatives with respect to $s$ are analytic at $s=0$. 
This implies, for instance, that the values $E_{L_2}(\tau, 0; 0)$ and $E_{L_2}'(\tau, 0; 0)$ are defined and finite,
and that we can expand $E_{L_2}(\tau, s;0)$ into its Taylor series expansion around $s=0$. 
By Proposition \ref{MSEis}, the Eisenstein series $E_{L_2}(\tau, 0; 0)$ has a meromorphic continuation 
$E_{L_2}^{\star}(\tau, s) = E_{L_2}^{\star}(\tau, s; 0) = \Lambda(s+1, \eta)E_{L_2}(\tau, s;0)$
to all $s \in {\bf{C}}$ which satisfies an even functional equation $E_{L_2}^{\star}(\tau, s) = E_{L_2}^{\star}(\tau, -s)$. 
Comparing the corresponding Taylor series expansions around $s = 0$ as we may, we then see
that for any $s \in {\bf{C}}$ with $0 \leq \Re(s) <1$ we have the relation 
\begin{align*} E_{L_2}^{\star}(\tau, 0) + E_{L_2}^{\star \prime}(\tau, 0) s + O(s^2) 
&= E_{L_2}^{\star}(\tau, 0) - E_{L_2}^{\star \prime}(\tau, 0) s + O(s^2), \end{align*}
equivalently 
\begin{align*} E_{L_2}^{\star \prime}(\tau, 0) s + O(s^2) &=  - E_{L_2}^{\star \prime}(\tau, 0) s + O(s^2). \end{align*}
Taking the limit as $\Re(s) \rightarrow 0$, we see that 
$E^{\star \prime}_{L_2}(\tau, 0)  = \Lambda(1, \eta) E'_{L_2}(\tau, 0; 0) + \Lambda'(1, \eta) E_{L_2}(\tau, 0;0)$ must vanish.
Now, if we multiply $(\ref{2.20})$ by $\Lambda(s+1, \eta)$, viewing this as a scalar in $\tau$ which commutes with $L_2$,
we get the corresponding relation of completed Eisenstein series 
\begin{align}\label{2.20*} L_2 E_{L_2}^{\star}(\tau, s; 2) &= \frac{1}{2}(s-1) E_{L_2}^{\star}(\tau, s: 0)\end{align}
Differentiating both sides of $(\ref{2.20*})$ in $s$ and evaluating at $s=0$, using that $E_{L_2}^{\star \prime}(\tau, 0; 0)=0$, we obtain 
\begin{align}\label{CFI} 2 L_2 E^{\star \prime}_{L_2}(\tau, 0; 2) = E_{L_2}^{\star}(\tau, 0; 0). \end{align}
On the other hand, since the weight-two Eisenstein series $E_{L_2}(\tau, s; 2)$ is incoherent, 
with odd functional equation and vanishing central value $E_{L_2}(\tau, 0; 2) = 0$ (Proposition \ref{incoherent}), we see that 
\begin{align*} E^{\star \prime}_{L_2}(\tau, 0;2) =\Lambda'(1,\eta)E_{L_2}(\tau, 0; 2) + \Lambda(1, \eta) E^{\prime}_{L_2}(\tau, 0; 2)
= \Lambda(1, \eta) E_{L_2}^{\prime}(\tau, 0; 2). \end{align*} Hence, the completed functional identity $(\ref{CFI})$ is equivalent to 
\begin{align*} 2 L_2 \Lambda(1, \eta) E^{\prime}_{L_2}(\tau, 0; 2) &= \Lambda(1, \eta) E_{L_2}(\tau, 0; 0). \end{align*}
Dividing out by the scalar $\Lambda(1, \eta)$ on each side, we obtain the claimed identity $2L_2 E'_{L_2}(\tau, 0; 2) = E_{L_2}(\tau, 0; 0)$. \end{proof}

\end{document}
